\chapter{Otosclerosis}
\index{Otosclerosis}

\medskip
\noindent\textit{This chapter is for general education. Gradual hearing loss or persistent tinnitus should be assessed.}

\section*{What it is}
Otosclerosis is abnormal bone growth around a small middle-ear bone, often the stapes. The stapes cannot move normally, so sound transmission is reduced. It commonly causes gradual conductive hearing loss and may affect one or both ears.

\section*{Symptoms and assessment}
Symptoms may include progressive hearing loss, tinnitus, dizziness, or balance problems. Assessment usually includes looking in the ear, hearing tests, and review by an audiologist or ear, nose, and throat specialist. A CT scan may be used when the diagnosis or surgical plan is uncertain.

\section*{Treatment}
Hearing aids can make sounds louder but do not remove the underlying change. Selected people may have stapes surgery, such as stapedotomy or stapedectomy, to improve sound transmission. The specialist should explain possible benefits, alternatives, and risks, including dizziness, taste changes, tinnitus, or worsening hearing.

\section*{Living with hearing loss}
Use communication strategies, hearing technology, and support services as needed. Reduce background noise, face the person speaking, and ask for information in accessible formats. Do not put objects or unprescribed drops into the ear.

\section*{When to Seek Medical Care}
Arrange assessment for gradually worsening hearing, persistent or troublesome tinnitus, dizziness, or balance problems. Seek urgent help for sudden hearing loss, facial weakness, severe vertigo, severe headache, or neurological symptoms.

\section*{Sources and further reading}
\begin{itemize}
\item NHS. \textit{Otosclerosis}. \url{https://www.nhs.uk/conditions/otosclerosis/}
\end{itemize}
